% docs/manual/operations/chapters/08-recovery.tex
% 第 8 章：故障恢复与排错

\chapter{故障恢复与排错}

长跑任务总会出故障：网络断、磁盘满、机器重启、OOM。本章详解每种故障下"最快回到能跑"的步骤。

\section{Ctrl-C 中断}

\subsection{基础}

\openbench{} 在每个 task 完成时立即写 cache + 输出。Ctrl-C 中断只丢失"正在进行"的那个 task。

\subsection{续跑}

\begin{minted}{bash}
openbench run config.yaml      # 自动从 cache 命中处续上
\end{minted}

如果你的 \yamlkey{project.force: false}（默认），增量缓存会跳过已完成 task。

\subsection{留下半成品}

正在写的 zarr / CSV / 图可能不完整。\openbench{} 没主动检测。如果重跑后某些输出缺失或报错，删那个 task 的 cache 再跑：

\begin{minted}{bash}
rm output/<run>/scratch/cache/<task_key>.json
rm -rf output/<run>/data/<task_key>.zarr
openbench run config.yaml
\end{minted}

\section{异常崩溃后续跑}

\subsection{看 traceback}

\begin{minted}{bash}
# 看 .log 末尾
tail -100 output/<run>/run.log

# 找 traceback
grep -A 20 "Traceback" output/<run>/run.log | tail -40
\end{minted}

定位失败的 task。

\subsection{根因分类}

\begin{longtable}{p{3cm} p{4cm} p{6cm}}
  \toprule
  根因 & 提示 & 修复 \\
  \midrule
  \endhead
  OOM & 进程被 kill；no traceback & 缩 \yamlkey{num\_cores} 或区域 \\
  磁盘满 & "No space left" & 清旧输出 \\
  HDF5 lock & "Resource temporarily" & 见 第 6 章 \\
  数据问题 & "KeyError"、坐标错 & 检查 NC 文件 \\
  代码 bug & 长 traceback & 提 GitHub Issue \\
  \bottomrule
\end{longtable}

\subsection{续跑}

修了根因后：

\begin{minted}{bash}
openbench run config.yaml      # cache 自动跳已完成
\end{minted}

\section{残留 lock 文件}

\subsection{HDF5 lock}

\openbench{} 已禁用 HDF5 lock；很少残留。如有：

\begin{minted}{bash}
find /data/refs -name "*.lock" -delete
find /tmp -name "*hdf5*" -mtime +1 -delete
\end{minted}

\subsection{zarr lock}

zarr 写入用文件级 atomic rename，不需要 lock。若发现 \texttt{*.partial} 或 \texttt{*.tmp}：

\begin{minted}{bash}
find output/<run>/data -name "*.partial" -delete
find output/<run>/data -name "*.tmp" -delete
\end{minted}

\section{磁盘满}

\subsection{识别哪占空间}

\begin{minted}{bash}
du -sh output/* | sort -h | tail
# 通常 data/（zarr）最大；其次 figures/
\end{minted}

\subsection{清理优先级}

\begin{enumerate}
  \item 旧 run 的 \texttt{scratch/} 与 \texttt{tmp/}（cache，可重建）
  \item 旧 run 的 \texttt{data/}（zarr，可重建但费时）
  \item 旧 run 的 \texttt{\_minted/}（LaTeX 中间）
  \item 旧 run 的 \texttt{figures/}（重新生成代价不大）
  \item 最后才动 \texttt{metrics/scores/comparisons/}（这些是结论数据）
\end{enumerate}

\subsection{预估磁盘需求}

\begin{itemize}
  \item Reference 数据：几十 GB - 几 TB（视集合）
  \item Sim 数据：依模型 + 时段
  \item Cache：sim + ref 的 0.5-2 倍
  \item 输出（不含 cache）：几十 MB - 几 GB
\end{itemize}

\textbf{规划}：可用磁盘 = (sim + ref) × 3。

\subsection{配额超限}

HPC 上 \texttt{quota} 超限会导致写失败。先：

\begin{minted}{bash}
quota -u $USER
df -h /scratch/$USER
\end{minted}

\section{Nested parallelism 死锁}

历史问题（2026-04-03 已修主流程，但自定义脚本可能复现）：

\subsection{症状}

\begin{itemize}
  \item 进程 CPU 100\% 但不出进度
  \item py-spy 看到 joblib 内部等锁
  \item 重启后第一个 task 也卡
\end{itemize}

\subsection{诊断}

\begin{minted}{bash}
# pip install py-spy
sudo py-spy dump --pid $(pgrep -f 'openbench run')
\end{minted}

如果调用栈深处反复出现 \texttt{joblib} + \texttt{ThreadPool} + \texttt{Lock}，是嵌套并行死锁。

\subsection{修复}

\begin{minted}{bash}
export OMP_NUM_THREADS=1
export MKL_NUM_THREADS=1
export OPENBLAS_NUM_THREADS=1
openbench run config.yaml
\end{minted}

详见第 4 章。

\section{partial 状态恢复}

\openbench{} 的 \cli{run} 在某些 task 失败、其他成功时返回 \texttt{partial}。

\subsection{识别}

\begin{minted}{text}
✗ Evaluation completed with errors
  Output: output/<run>
  - [evaluation] Latent_Heat sim=X ref=Y: ...
  - [evaluation] Net_Radiation sim=X ref=Z: ...
\end{minted}

\subsection{重跑只失败的}

最简单：直接 \cli{openbench run config.yaml}，cache 让成功的 task 跳过，失败的会重试（可能再失败如果根因没解决）。

\subsection{强制全量重算}

\begin{minted}{bash}
openbench run config.yaml --cores N   # CLI 临时覆盖
# 或
yaml: project.force: true
\end{minted}

\section{从 cache 损坏恢复}

详见第 5 章 "cache 损坏检测与修复" 部分。简言之：删损坏的 zarr + 对应 task cache，重跑。

\section{升级 \openbench{} 后旧输出}

升级后某些 metric 实现可能变。\textbf{当前 v3.0a1 cache key 不含版本}，所以旧 cache 会被误命中。

\subsection{安全做法}

\begin{minted}{bash}
# 升级前
pip show openbench | grep Version

# 升级
pip install -U openbench

# 升级后第一次 run
openbench run config.yaml --cores N   # 但 cache 命中可能用旧版数值
# 或显式
yaml: project.force: true   # 全量重算保险
\end{minted}

或干脆：

\begin{minted}{bash}
rm -rf output/<old-run>/scratch
rm -rf output/<old-run>/data
openbench run config.yaml
\end{minted}

\section{不可恢复的失败：重启评估}

最后手段：彻底重新开始。

\begin{minted}{bash}
# 1. 备份关键产物（如有）
cp -r output/<run>/figures /tmp/saved-figures
cp output/<run>/run.log /tmp/saved-run.log

# 2. 删整个 run 目录
rm -rf output/<run>

# 3. 重跑
openbench run config.yaml
\end{minted}

\section{诊断 checklist 速查}

按出现频率排序：

\begin{enumerate}
  \item 看 \texttt{output/<run>/run.log} 末尾 100 行 \texttt{tail -100 ...}
  \item 找最后一个 ERROR \texttt{grep ERROR ...}
  \item 看进程是否还活着 \texttt{pgrep -af openbench}
  \item 看磁盘是否满 \texttt{df -h}
  \item 看内存是否够 \texttt{free -g}
  \item 看 SLURM 状态（如适用） \texttt{squeue -u \$USER}
  \item 看是否有 zombie HDF5 lock \texttt{find /data -name "*.lock"}
\end{enumerate}

\section{完结}

到这里你已读完运维卷全部 8 章，覆盖部署 → 安装 → 远程 → 性能 → 缓存 → 站点 → 监控 → 恢复。下一步：

\begin{itemize}
  \item 看附录 A-E 做参考查阅
  \item 部署到生产前过一遍第 4 章"调优 checklist"
  \item 大评估前用 SLURM + 第 7 章监控建议
  \item 出问题先按本章诊断 checklist
\end{itemize}

谢谢使用 \openbench{} —— 也欢迎贡献你的运维经验回主仓库。
